% DEIXIS evidence report, IEEEtran export for XeLaTeX.
% Compile with: latexmk -xelatex report-synthetic-report-v12.tex
% Needs XeLaTeX, BibTeX and the IEEEtran class and IEEEtran.bst style; they are not included.
% Export notes: none

\documentclass[journal]{IEEEtran}
\usepackage{fontspec}
\usepackage{amsmath}
\usepackage{amssymb}
\usepackage{bm}
\usepackage{xcolor}
\usepackage{cite}
\usepackage{url}
\usepackage{booktabs}
\usepackage{longtable}
\usepackage{array}

\begin{document}

\renewcommand{\abstractname}{Özet}
\renewcommand{\IEEEkeywordsname}{Dizin Terimleri}
\renewcommand{\refname}{Kaynaklar}

\title{SYNTHETIC report \& \ensuremath{\alpha}}

\maketitle

\noindent Kanıt raporu \ensuremath{\cdot} V12

\section*{I. Giriş}

SYNTHETIC claim. \cite{Synth26}

\section*{Rapor notları}

DEIXIS ile üretildi.

Atıf çapaları ilgili pasajlarda veya hücrelerde bulundu. Raporu yazan model, ek bir inceleme çağrısında 1 bölümün 1 tanesindeki iddiaları atıf yapılan pasaj ve hücrelerle karşılaştırıp 1 olası sorun işaretledi. Bu, modelin okumasıdır; hakem incelemesi değildir ve hataları kaçırabilir. Kod, her pasajın ilgili iddiayı destekleyip desteklemediğini denetlemedi. Model, yeniden yazılan 2 cümlenin kaynaklarıyla artık uyuşmayabileceğini işaretledi; cümleler özgün hâline döndürüldü. İnceleme modelin temel sürümünü kapsar; insan düzenlemeleri incelenmedi.

\noindent\textbf{Model bulguları}

\begin{itemize}
\item I. Giriş \ensuremath{\cdot} Sayı \ensuremath{\cdot} SYNTHETIC finding 1: a \ensuremath{\geq} b \& c.
\end{itemize}

Bu çıktı raporun metnini ve kaynakçasını taşır. Alıntı pasajları, tablo hücresi kanıtları ve konumlanmış atıf çapaları DEIXIS'te kalır.

\bibliographystyle{IEEEtran}

\bibliography{report-synthetic-report-v12}

\end{document}
